\documentclass[10pt,conference]{IEEEtran}
\usepackage{amsmath,amssymb,booktabs,array}
\usepackage[hidelinks]{hyperref}

\title{Fixed-Graph Vanilla and Relay Belief-Propagation Decoders in Chisel}
\author{Anonymous Authors}

\begin{document}
\maketitle

\begin{abstract}
This manuscript records the current design, implementation, and validation of
\texttt{chipsLDPC}, a Chisel/CIRCT generator for fixed-Tanner-graph quantum
low-density parity-check decoders.  One shared statically wired min-sum core is
controlled by two autonomous decoders: early-terminating vanilla BP and
Relay-BP-S with persistent marginals, deterministic pseudorandom per-variable memory
coefficients, zero-bubble leg transitions, and minimum-prior-score solution
selection.  Every flooding iteration uses one registered check-node clock and
one registered variable-node clock.  The maintained experiments specialize
both decoders to the same reduced Z-check detector graph of the
$[[144,12,12]]$ bivariate-bicycle memory circuit, compare the exact emitted
SystemVerilog against an independent Scala golden model, execute deterministic
shots in parallel Verilator shards, and report logical word failure and
measured architecture clocks.  A runtime-budget meta-experiment reuses one
maximum-capacity artifact and the same shot corpus to plot failure probability
against mean executed iterations.  This document states the arithmetic,
controller state, data-structure delta, cycle equations, simulator checks, and
file-to-file workflow explicitly so that each reported quantity is traceable to
an implementation symbol.  FPGA synthesis, placement, frequency, resource,
power, and board-execution claims remain outside the current artifact.
\end{abstract}

\section{Research Objective and Current Boundary}
Quantum LDPC decoding research must permit rapid algorithmic changes while
keeping the generated hardware, arithmetic, and timing boundaries inspectable.
The purpose of \texttt{chipsLDPC} is therefore a small collection of explicit
generators rather than a general decoder framework.  A parity-check matrix is
elaboration-time data.  Chisel constructs a fixed circuit for that graph,
CIRCT emits inspectable FIRRTL-dialect MLIR and synthesizable SystemVerilog, and
Verilator checks both generated designs and an exact emitted artifact.

The implemented decoder milestone contains the following:
\begin{itemize}
  \item reusable combinational CNU and VNU kernels;
  \item Valls-style constant shift-add scaling and Relay-style iteration-ramp
        scaling as elaboration-time CNU policies;
  \item an FPGA-paper-style reduced Relay multiplier and deterministic
        per-variable LFSR coefficient source;
  \item an elaboration-only Tanner graph and local node-port plan;
  \item a full static network with one CNU per row, one VNU per column, one
        fixed connection per nonzero of $H$, and a two-phase local scheduler;
  \item an immediate-commit view, a registered result view, and a combinational
        residual/convergence checker;
  \item autonomous vanilla and Relay-BP controllers with optional per-frame
        runtime limits and backpressured outputs;
  \item two runnable, exact-RTL BB144 experiments sharing the same physical
        samples, logical observables, sharding, and reporting path;
  \item a runtime iteration-budget sweep over either emitted decoder; and
  \item a separate binary Gauss--Jordan component family for future
        post-processing experiments.
\end{itemize}
The shared static core instantiates the base min-sum VNU once.  Vanilla BP
drives its bias input from the immutable prior; Relay BP computes a new bias
from the immutable prior, the registered previous marginal, and a coefficient
that is stable for one leg.  Thus Relay support is integrated without a second
Tanner network or a duplicated VNU datapath.  The emitted benchmark artifacts
are single-frame engines: they do not bank independent decoder states or overlap
shots.  Host interfaces, FPGA implementation, and board execution remain
outside the implemented boundary.  The independent Gauss--Jordan and filtered
OSD-0 subsystems remain available but are not instantiated in either BP-only
experiment artifact described here.

\section{Algorithmic Background}
Let $H\in\{0,1\}^{m\times n}$ be a parity-check matrix.  Check $i$ is adjacent
to variables $\mathcal{N}(i)$ and variable $j$ is adjacent to checks
$\mathcal{M}(j)$.  The incoming variable-to-check message is $R_{i,j}$, the
outgoing check-to-variable message is $\sigma_{i,j}$, and $S_i$ is a syndrome
bit.  The syndrome-based min-sum check update is
\begin{equation}
 \sigma_{i,j}^{(t)} =
 (-1)^{S_i\oplus
 \bigoplus_{j'\in\mathcal{N}(i)\setminus j}
 [R_{i,j'}^{(t-1)}<0]}
 g_t\!\left(
 \min_{j'\in\mathcal{N}(i)\setminus j}|R_{i,j'}^{(t-1)}|
 \right).
 \label{eq:cnu}
\end{equation}
The variable update and marginal are
\begin{align}
 R_{i,j}^{(t)} &= B_j^{(t)}+
 \sum_{i'\in\mathcal{M}(j)\setminus i}\sigma_{i',j}^{(t)},
 \label{eq:vnu}\\
 M_j^{(t)} &= B_j^{(t)}+
 \sum_{i'\in\mathcal{M}(j)}\sigma_{i',j}^{(t)}.
 \label{eq:marginal}
\end{align}
The hard decision is deliberately strict:
\begin{equation}
 \hat e_j^{(t)}=\operatorname{HD}(M_j^{(t)})=
 \begin{cases}
  1,&M_j^{(t)}<0,\\
  0,&M_j^{(t)}\geq 0.
 \end{cases}
 \label{eq:hard-decision}
\end{equation}
Thus a zero marginal means no correction.  The implementation takes this
decision from the unsaturated working-width total and then saturates the stored
marginal.  Vanilla BP uses
\begin{equation}
 B_j^{(t)}=\Lambda_j,
 \qquad
 R_{i,j}^{(0)}=+\Lambda_j,
 \label{eq:vanilla-bias}
\end{equation}
where the benchmark prior $\Lambda_j$ is an unsigned quantized magnitude.
For the ramp policy used by both maintained experiments, the corresponding
real-valued scale factor is
\begin{equation}
 \alpha_t=
 \begin{cases}
  1-2^{-t},&1\leq t\leq w_m,\\
  1,&t>w_m,
 \end{cases}
 \label{eq:ramp-alpha}
\end{equation}
The exact unsigned-integer operation used in Eq.~\eqref{eq:cnu} is
\begin{equation}
 g_t(x)=
 \begin{cases}
  x-\left\lfloor x/2^t\right\rfloor,&1\leq t\leq w_m,\\
  x,&t>w_m,
 \end{cases}
 \label{eq:ramp-integer}
\end{equation}
which is also $x-(x\mathbin{\texttt{>>}}t)$ in the Chisel implementation.

Relay BP replaces the static VNU bias by a memory mixture.  For Relay leg
$\ell$, local iteration $t$, coefficient
$\bar\beta_{\ell,j}=b_{\ell,j}/2^f$, and
$\gamma_{\ell,j}=1-\bar\beta_{\ell,j}$, the ideal equation is
\begin{equation}
 B_j^{(\ell,t)}=
 \bar\beta_{\ell,j}\Lambda_j+
 (1-\bar\beta_{\ell,j})M_j^{(\ell,t-1)}.
 \label{eq:memory}
\end{equation}
The hardware deliberately uses a reduced partial-product calculation.  Define
$x_r^+$ as bit $r$ of the unsigned magnitude $|x|$ and define the restored sign
$\epsilon(x)$ by
\begin{align}
 |x|&=\sum_{r=0}^{w_a-1}x_r^+2^r,\\
 \epsilon(x)&=\begin{cases}-1,&x<0,\\+1,&x\geq0,\end{cases}\\
 \rho_f(x,b)&=\epsilon(x)\sum_{r=0}^{w_a-1}x_r^+
 \left\lfloor\frac{b2^r}{2^f}\right\rfloor.
 \label{eq:reduced-product}
\end{align}
The implemented Relay bias is then
\begin{equation}
 \begin{split}
 B_j^{(\ell,t)}=\operatorname{sat}_{w_a}\big(&
 \rho_f(\Lambda_j,b_{\ell,j})+M_j^{(\ell,t-1)}\\
 &-\rho_f(M_j^{(\ell,t-1)},b_{\ell,j})\big),
 \end{split}
 \label{eq:relay-bias}
\end{equation}
where every selected partial product is right-shift-truncated before the sum;
this is not a full product followed by one truncation.  Here
\begin{equation}
 \operatorname{sat}_{w}(x)=
 \min\!\left(2^{w-1}-1,\max(-2^{w-1},x)\right).
 \label{eq:saturation}
\end{equation}
The first state after load has $M_j^{(0,0)}=\Lambda_j$.  At a Relay leg
boundary the newly computed saturated marginal is retained while every
$R_{i,j}$ is synchronously restored to $+\Lambda_j$.

Convergence is a separate binary computation.  For hard-decision vector
$\hat e$, the residual is
\begin{equation}
 r_i^{(t)}=S_i\oplus\bigoplus_{j:H_{i,j}=1}\hat e_j^{(t)},
 \qquad \text{converged}\iff \left(\bigvee_i r_i^{(t)}\right)=0.
 \label{eq:residual}
\end{equation}

\section{Programming Model and Reading Order}
\subsection{Three distinct representations}
The code uses Scala, Chisel, and emitted hardware for different purposes.  A
plain Scala \texttt{Seq}, \texttt{Vector}, case class, loop, or pattern match
runs once while elaborating a particular circuit.  A Chisel \texttt{Vec},
\texttt{UInt}, \texttt{SInt}, \texttt{Bool}, \texttt{Mux}, or \texttt{when}
describes hardware that exists at runtime.  FIRRTL-dialect MLIR and
SystemVerilog are generated artifacts, not additional hand-maintained decoder
implementations.  Keeping these levels distinct is the key to reading the
repository.

For example, \texttt{nodes.checks.map} in
\path{StaticTannerDatapath.scala} is a Scala traversal that creates one module
per check during elaboration.  It is not a runtime iterator.  Conversely,
\texttt{io.step.fire} is a hardware condition evaluated on every clock cycle.
The Scala match over \texttt{CheckScale} selects hardware at elaboration time;
the Chisel \texttt{MuxLookup} in the ramp policy selects a constant-shift result
at runtime.

\subsection{Module conventions}
Combinational tops are \texttt{RawModule}s, so they acquire no accidental
clock or reset.  Stateful components are \texttt{Module}s and place registers
at intentional phase or protocol boundaries.  \texttt{RegInit} is used for
control or externally required reset state; \texttt{RegEnable} and unreset
\texttt{Reg} are used for wide datapath state whose validity is guaranteed by
the load/step protocol.  This reset policy avoids manufacturing a large reset
network for data that is overwritten before use.

Classes are \texttt{final} unless extension is part of the design.  Immutable
case classes carry configuration or reference values.  Composition and
elaboration-time policy objects are preferred to hardware inheritance.  The
package-restricted graph records prevent callers from constructing inconsistent
local port maps while still exposing the validated result to decoder code.

\subsection{Recommended source traversal}
Table~\ref{tab:reading-order} gives the shortest path through the maintained
implementation.  A reader following it encounters every contract before its
first important use.

\begin{table*}[t]
\caption{Source-oriented reading order.}
\label{tab:reading-order}
\centering
\scriptsize
\begin{tabular}{@{}p{0.035\textwidth}p{0.255\textwidth}p{0.33\textwidth}p{0.31\textwidth}@{}}
\toprule
Step & Source & Read for & Then inspect \\
\midrule
1 & \path{Config.scala} & Width, degree, scaling, and Relay defaults; all are immutable Scala values. & The \texttt{require} clauses, then the locked \texttt{RelayDefaults}. \\
2 & \path{graph/TannerGraph.scala} & Canonical sparse $H$, row-major global edges, and derived column adjacency. & \texttt{fromRows}; note which checks happen before hardware elaboration. \\
3 & \path{graph/TannerNodeGraphs.scala} & Translation from each global edge to local CNU and VNU ports. & \texttt{connections}; distinguish node ID from local port number. \\
4 & \path{Protocol.scala} & Hardware bundles crossing arithmetic and phase boundaries. & \texttt{CheckResult} versus the wire-view \texttt{CheckMessage}. \\
5 & \path{Arithmetic.scala} & Canonical sign, conversion, balanced sums, explicit widths, clipping, and scale policies. & Operators that expand or retain width and the elaboration/runtime split in \texttt{scale}. \\
6 & \path{TwoMin.scala} and \path{CheckNode.scala} & Balanced first/second minimum, shared parity, scaling, and compressed broadcast metadata. & Why equality selects the second minimum and why duplicates remain valid. \\
7 & \path{VariableNode.scala} & Selection, sign--magnitude conversion, one shared total, and total-minus-own messages. & \texttt{VariableLogic.apply} and \texttt{VariableLogic.connect}. \\
8 & \path{ConvergenceChecker.scala} & Combinational $H\hat e\oplus S$ specialized from the Scala graph. & Native XOR and OR reductions. \\
9 & \path{StaticTannerDatapath.scala} & Full graph elaboration, load initialization, static edge wiring, immediate commit, registered result, and message restart. & \texttt{StaticVariableStage}, then \texttt{StaticTannerCore}. \\
10 & \path{BpDecoder.scala} & Common result bundle and the vanilla BP state machine. & \texttt{VanillaBpConfig}, \texttt{BpDecoderResult}, then \texttt{VanillaBpDecoder}. \\
11 & \path{RelayBias.scala}, \path{RelayCoefficients.scala}, and \path{RelayBpDecoder.scala} & Reduced products, LFSR coefficients, Relay state, leg transitions, and solution selection. & \texttt{RelayBias.reducedMultiply}, \texttt{LfsrRelayCoefficients}, then \texttt{RelayBpDecoder}. \\
12 & \path{Reference.scala} & Independent integer equations and complete vanilla/Relay runs. & \texttt{runVanilla}, \texttt{runRelay}, and \texttt{lfsrBeta}. \\
13 & \path{BpOnlyExperiment.scala} and \path{RelayBpExperiment.scala} & Packed artifact wrappers, golden records, runtime limits, and RTL emission. & The artifact class before the executable object. \\
14 & \path{benchmarks/bb144/sim_main.cpp} & Exact-golden and many-shot modes, sparse reconstruction, timing assertions, and logical classification. & \texttt{decode}, \texttt{exactMain}, then \texttt{benchmarkMain}. \\
15 & \path{examples/} and \path{benchmarks/rtl_shards/} & End-to-end commands, hashed shard manifests, merge rules, and reports. & Each \texttt{run.sh}, then the adapter, workflow, and report. \\
16 & \path{Generate.scala} and \path{GaussJordan/} & Inspection tops and independent one-bit post-processing PEs. & Treat these as emission support and a separate subsystem. \\
\bottomrule
\end{tabular}
\end{table*}

\section{Configuration and Static Graph Elaboration}
\subsection{Validated configuration objects}
\texttt{Quantization} stores edge magnitude width $w_m$ and accumulator width
$w_a$, requiring $w_a\geq w_m+1$.  \texttt{CheckConfig} adds check degree and
a \texttt{CheckScale}; checks require degree at least two.
\texttt{VariableConfig} accepts positive degree and derives a working width.
\texttt{RelayFormat} describes an unsigned fixed-point integer coefficient
$b$ with $f$ fractional bits and requires \texttt{betaBits}$>f$.
The two maintained experiment profiles are deliberately distinct:
\begin{align}
 (w_m,w_a)_{\mathrm{vanilla}}&=(4,7),
 &\texttt{scale}&=\texttt{RampScale}(4),\label{eq:vanilla-format}\\
 (w_m,w_a)_{\mathrm{Relay}}&=(4,5),
 &(w_b,f,M)&=(4,3,8).\label{eq:relay-format}
\end{align}
Here $M=2^f=8$ and $b/M$ represents $\bar\beta$.  The Python problem source
uses the separate prior-quantization scale $S_\Lambda=2$.  These values are
defined by \texttt{RelayDefaults.q}, \texttt{paperQ}, \texttt{paperFormat},
\texttt{priorScale}, and \texttt{memoryScale} in \path{Config.scala}.  Keeping
the seven-bit vanilla accumulator separate prevents the established vanilla
experiment from changing when the five-bit Relay paper profile is selected.

\subsection{One canonical graph source}
The factory \texttt{TannerGraph.fromRows} is the only constructor for the
static parity-check graph.  Its arguments are \texttt{variableCount} and
\texttt{rowOnes}.  For $H\in\{0,1\}^{m\times n}$, \texttt{rowOnes(i)} lists
the variable indices in $\mathcal{N}(i)$.  The
constructor sorts every row, rejects a nonpositive variable count, an empty
check set, empty rows, duplicate or out-of-range entries, and any variable with
no incident check.  It then assigns edge IDs in increasing check order and
increasing variable order.  Column adjacency, check-edge IDs, and variable-edge
IDs are derived rather than accepted as redundant inputs.

This object contains no Chisel types.  It adds no ports, memories, registers,
or runtime router.  Graph validation and transpose construction finish before
the first hardware module is created.

\texttt{TannerNodeGraphs.from} derives rooted node views without renumbering an
edge.  A connection record
\begin{equation}
 (e,i,p_i,j,p_j)
\end{equation}
contains global \texttt{edgeId} $e$, check \texttt{checkId} $i$, local
\texttt{checkPort} $p_i$, variable \texttt{variableId} $j$, and local
\texttt{variablePort} $p_j$.  Node IDs select a generated module; port numbers
select a position inside that module's degree-dependent \texttt{Vec}.  Both are
needed because the same global edge generally has different local positions at
its two endpoints.  The constructor asserts that the connection sequence still
contains exactly the canonical global edge IDs.

\subsection{Steane elaboration example}
The concrete graph-level example is one CSS sector of the Steane
$[[7,1,3]]$ code~\cite{steane1996}.  With zero-based variable indices,
\begin{equation}
 H_{\mathrm{Steane}} =
 \begin{bmatrix}
 0&0&0&1&1&1&1\\
 0&1&1&0&0&1&1\\
 1&0&1&0&1&0&1
 \end{bmatrix}.
\end{equation}
Consequently,
\begin{equation}
 \texttt{rowOnes}=
 [[3,4,5,6],[1,2,5,6],[0,2,4,6]].
\end{equation}
The graph has three degree-four checks, seven variables with degrees
$(1,1,2,1,2,2,3)$, and twelve edges.  Its identical Hamming-code blocks
$H_X=H_Z$ let the same topology represent either CSS sector independently.
The exact rooted views and all twelve five-field connection records are locked
by \path{TannerNodeGraphsSpec.scala}.

\section{Numeric and Interface Contract}
Table~\ref{tab:numeric-contract} is normative for the current implementation.

\begin{table}[t]
\caption{Implemented fixed-width semantics.}
\label{tab:numeric-contract}
\centering
\small
\begin{tabular}{@{}p{0.30\columnwidth}p{0.63\columnwidth}@{}}
\toprule
Property & Contract \\
\midrule
Edge representation & Separate sign and unsigned magnitude. \\
Complete zero & A zero-magnitude \texttt{SignMag} is nonnegative. \\
CNU metadata & Sign is parity metadata until paired with the selected magnitude. \\
Minimum tie & Equal values occupy both minimum positions. \\
Selector & Use second iff the incoming magnitude equals the unscaled first minimum. \\
Decision tie & Marginal zero produces correction zero. \\
Accumulator & Signed two's-complement with an explicit guard width. \\
Saturation & Only at named stored or edge interfaces. \\
Initialization & Static V2C edges start as positive copies of their variable prior. \\
\bottomrule
\end{tabular}
\end{table}

The complete edge message \texttt{SignMag} canonicalizes zero through
\texttt{Arithmetic.canonicalSign}.  The CNU's \texttt{EdgeMeta.sign} is not by
itself a complete sign--magnitude number: it is exclusive-parity metadata that
may be one even when a subsequently selected magnitude is zero.  The VNU
combines the fields and canonicalizes zero during conversion.  This distinction
preserves a clean split between sign parity and magnitude selection.

The VNU working width is
\begin{equation}
 w_{\mathrm{work}}=w_a+
 \max\left(1,\left\lceil\log_2(d_v+1)\right\rceil\right).
\end{equation}
Every prior and incoming message is padded to this width.  The balanced sum
uses explicit fixed-width addition; the result is clipped once into the stored
$w_a$ marginal.  Each total-minus-own result is separately converted and
clipped into an unsigned $w_m$ magnitude.  In Chisel, padding, signed casts,
fixed-width arithmetic, and expanding arithmetic are written explicitly, so a
reader should not infer mathematical unbounded integers from a Scala-looking
expression.

\subsection{Bundle-level broadcast protocol}
\texttt{Protocol.scala} makes the compression boundary visible:
\begin{itemize}
  \item \texttt{MinPair} carries the shared first and second magnitudes;
  \item \texttt{EdgeMeta} carries one sign and one \texttt{useSecond} bit;
  \item \texttt{CheckResult} stores one \texttt{MinPair} and a vector of edge
        metadata, not one copied pair per edge;
  \item \texttt{CheckMessage} is the VNU-side wire view of one broadcast pair
        plus one edge's metadata; and
  \item \texttt{VariableResult} carries the stored marginal, hard decision,
        and one outgoing sign--magnitude message per incident edge.
\end{itemize}
The static graph therefore registers $2w_m$ broadcast magnitude bits per CNU,
plus two metadata bits per edge.  It does not register an edge-expanded array
of duplicated check-to-variable magnitudes.  Physical fanout and placement are
left to synthesis and implementation tools; the source-level architecture does
not claim a routing result before those tools run.

\section{Check-Node Datapath}
\subsection{Canonical exclusive sign}
\texttt{CheckNode} first turns every input into a canonical sign bit and forms
one XOR reduction with the syndrome.  The exclusive sign for port $j$ is the
global parity XOR the canonical sign of port $j$.  This shares a parity tree
across all outputs and implements Eq.~\eqref{eq:cnu} without building an
independent XOR tree per edge.

\subsection{Balanced duplicate-preserving two-minimum finder}
\texttt{TwoMinLogic} pads a degree $d_c$ input sequence to the next power of two
$P$ with the maximum representable magnitude.  Each tournament candidate
stores its winner and the direct losers encountered on that winner's path.  A
balanced, stable winner tree uses less-than-or-equal, so the left value wins a
tie and the equal right value remains on the loser path.  The first minimum is
the final winner; a balanced reduction over its direct losers produces the
second minimum.  Thus $(2,2,5)\mapsto(2,2)$ rather than dropping the duplicate.

The tournament uses $P-1$ comparisons and the loser-path reduction uses
$\log_2P-1$, for approximately $P+\log_2P-2$ comparisons.  This describes the
generated structure, not an achieved FPGA delay.  The separate \texttt{TwoMin}
\texttt{RawModule} is only a thin IO wrapper around the reusable logic object.

Each input whose magnitude equals the unscaled first minimum sets
\texttt{useSecond}.  If the first minimum is duplicated, all tied edges select
the second value, which is equal and therefore still correct.  No minimum index
or one-hot vector is routed.

\subsection{Scaling policies}
Three Scala policy values share the same CNU:
\begin{itemize}
  \item \texttt{NoScale} returns the minimum unchanged;
  \item \texttt{VallsScale(a,b)} implements
        $\alpha x=(x\gg a)+(x\gg b)$ with one expanding addition and unsigned
        clipping; and
  \item \texttt{RampScale(k)} implements $x-(x\gg t)$ for the elaborated
        choices $0\leq t\leq k$, with an unscaled default for larger encoded
        controls.
\end{itemize}
Only the ramp policy elaborates the leaf \texttt{scaleShift} port; Scala
\texttt{Option} and pattern matching remove it from other leaf interfaces.
The two shared minima are each scaled once before broadcast, avoiding a scale
unit on every edge.

The resulting CNU combinational path is conceptually
\begin{equation}
 \text{V2C state}\rightarrow
 \{\text{XOR},\text{two-min},\text{scale}\}\rightarrow
 \text{CNU result register}.
\end{equation}

\section{Variable-Node Datapaths}
\subsection{Shared base logic}
\texttt{VariableLogic} is the common combinational implementation.  Each input
uses \texttt{useSecond} to select a broadcast magnitude, combines it with the
edge sign, and converts the complete sign--magnitude value to a sign-extended
two's-complement value.  A balanced adder tree computes
$M_j=\Lambda_j+\sum_i\sigma_{i,j}$ once.  Each outgoing message is then
\begin{equation}
 R_{i,j}=M_j-\sigma_{i,j}.
\end{equation}
This total-minus-own construction replaces $d_v$ independent adder trees with
one shared sum and $d_v$ subtractors.  \texttt{VariableLogic.connect} is a
small wiring helper; \texttt{VariableNode} is a combinational \texttt{RawModule}
that exposes the logic through a stable bundle interface.

The base leaf accepts a signed $w_a$-bit prior, so leaf tests can cover positive
and negative values.  The current static graph's load interface intentionally
accepts an unsigned $w_m$-bit LLR magnitude and zero-extends it, matching the
present independent-error examples with $p<1/2$.  There is no hidden hard-coded
$+1$ prior in the generator.

\subsection{Relay bias as a core input}
Relay arithmetic is composed around the base VNU rather than implemented as a
second node class.  \texttt{RelayBias.reducedMultiply} in
\path{RelayBias.scala} implements Eq.~\eqref{eq:reduced-product};
\texttt{RelayBias.apply} forms Eq.~\eqref{eq:relay-bias} and performs the final
five-bit signed saturation.  \texttt{RelayBpDecoder} applies this function once
per variable and drives the resulting vector into \texttt{core.io.bias}.
Consequently \texttt{VariableLogic} and every Tanner-edge bundle are identical
between vanilla and Relay BP.

The base VNU combinational path is
\begin{equation}
 \begin{aligned}
 \text{CNU registers}&\rightarrow\text{select/convert}
 \rightarrow\text{balanced sum}\\
 &\rightarrow\text{subtract/clip}
 \rightarrow\text{VNU state register}.
 \end{aligned}
\end{equation}

\section{Convergence Checker}
\texttt{ConvergenceChecker} accepts the pure Scala \texttt{TannerGraph} in its
constructor.  Elaboration specializes one native XOR reduction per row of
$H$.  At runtime it XORs each parity with the loaded syndrome to produce the
residual in Eq.~\eqref{eq:residual}; one OR reduction followed by negation
produces \texttt{converged}.  The block is a combinational \texttt{RawModule}:
it owns neither state nor a stop policy.  In \texttt{StaticTannerCore} it
observes the combinational next VNU decisions and the registered syndrome.
The core presents the checker in
two forms: \texttt{commit} evaluates the combinational next VNU decisions on
the terminal VNU cycle, whereas \texttt{result} presents the same information
from registers on the following cycle.  Autonomous controllers use
\texttt{commit} to decide without adding a control bubble; the manually stepped
compatibility shell exports \texttt{result}.

\section{Graph-Generated Static Datapath}
\subsection{Constructor and hierarchy}
\texttt{StaticTannerCore(nodes,q,scale)} receives a fully validated
\texttt{TannerNodeGraphs} plan; it does not recompute graph topology internally.
The default arguments are \texttt{Quantization(4,7)} and
\texttt{RampScale(4)}.  It rejects any check degree below two when elaborating
the min-sum network.  \texttt{StaticTannerDatapath} is a compatibility shell
that ties \texttt{bias=prior}, disables message restart, and exposes the older
manual load/step/result interface.

Scala \texttt{map} calls instantiate one \texttt{StaticCheckStage} for every
check and one \texttt{StaticVariableStage} for every variable.  Suggested names
\texttt{cnu\_i} and \texttt{vnu\_j} keep the emitted hierarchy inspectable.
One traversal of \texttt{nodes.connections} performs all graph wiring:
\begin{itemize}
  \item VNU extrinsic state at \texttt{variablePort} drives CNU input at
        \texttt{checkPort};
  \item the CNU's single registered minimum pair fans out to each adjacent
        VNU wire view; and
  \item that edge's registered sign and selector accompany the pair.
\end{itemize}
There is no runtime crossbar, adjacency memory, node-ID decoder, or graph
configuration bus.  A different $H$ produces a different specialized artifact.

\subsection{Core protocol}
The shared core IO contains:
\begin{itemize}
  \item a flipped \texttt{Decoupled} \texttt{load} channel carrying $m$
        syndrome bits and $n$ unsigned priors;
  \item a flipped \texttt{Decoupled} \texttt{step} channel carrying the scale
        control; and
  \item a signed bias vector and a \texttt{restartMessages} control;
  \item the registered immutable prior vector;
  \item an immediate \texttt{commit} carrying the next $n$ marginals and
        decisions, $m$ residual bits, and \texttt{converged}; and
  \item a registered \texttt{result} with the same payload.
\end{itemize}
A \texttt{load.fire} event stores every syndrome and prior, sets each initial
marginal to the prior, sets every initial decision to false, and initializes
every variable-to-check edge to positive sign and that variable's prior
magnitude.  Reloading is supported when the VNU phase is not active.  On an
enabled VNU edge, \texttt{restartMessages=false} stores the computed extrinsic
messages.  If it is true, the stage still stores the newly computed marginal
and decision but replaces every extrinsic message by the positive immutable
prior.  This selective assignment in \texttt{StaticVariableStage} is the entire
datapath mechanism required for a Relay leg transition.

The top gives load priority: \texttt{step.ready} requires a prior successful
load, no active VNU phase, and no asserted \texttt{load.valid}.  Therefore a
step cannot consume uninitialized datapath state, and an asserted load cannot
silently race a step.  Both result views use \texttt{Valid}, not
\texttt{Decoupled}; backpressure is handled only at the enclosing autonomous
decoder's frame output.

\subsection{Two-cycle iteration}
Only three control registers are reset: \texttt{loaded}, \texttt{vnuPhase}, and
\texttt{resultValid}.  The cycle sequence after a completed load is:
\begin{enumerate}
  \item On a \texttt{step.fire} edge, every CNU stage captures its result from
        the current V2C state.  \texttt{result.valid} remains low.
  \item During the next cycle, \texttt{commit.valid} is high and its payload is
        the next VNU result computed from the CNU registers.  On the closing
        edge, every VNU stage captures that result.  The controller may also
        stop, advance a Relay leg, or launch sparse output from this commit.
  \item \texttt{result.valid} is high in the following cycle and presents the
        registered marginals, decisions, residual, and convergence flag.
  \item During that result cycle, \texttt{step.ready} is high again, so a new
        CNU phase may be accepted at the next edge.
\end{enumerate}
The node-update latency and initiation interval are therefore two cycles per
iteration.  Wide syndrome, CNU-result, prior, and VNU-state registers are not
reset.  Their enabling protocol guarantees they are written before use; this
is intentional and is tested at the external interface.  Controller-visible BP
latency is therefore exactly
\begin{equation}
 C_{\mathrm{BP}}=2I,
 \label{eq:bp-clocks}
\end{equation}
for $I$ executed iterations; there is no extra iteration-control clock.

\section{Autonomous Decoder Implementations}
\subsection{Common input and output records}
\texttt{StaticDecoderInput} in \path{StaticTannerDatapath.scala} carries the
runtime syndrome and unsigned prior vector.  \texttt{BpDecoderResult} in
\path{BpDecoder.scala} is shared by both autonomous controllers and contains
\texttt{success}, selected \texttt{correction}, final \texttt{marginal},
reported \texttt{residual}, \texttt{totalIterations}, \texttt{legsExecuted},
and \texttt{solutionsFound}.  Both controllers use a three-state
\texttt{sIdle}/\texttt{sRun}/\texttt{sOutput} machine and a
\texttt{Decoupled} frame output.  Once in \texttt{sOutput}, the registered
result remains stable until \texttt{io.out.fire}; output backpressure therefore
cannot modify decoder state.

\subsection{Vanilla BP controller}
\texttt{VanillaBpConfig} stores the node plan, elaborated iteration capacity
$T_{\max}$, quantization, and scale policy.  Its counter width is
\begin{equation}
 w_I=\max\!\left(1,\left\lceil\log_2(T_{\max}+1)\right\rceil\right).
 \label{eq:vanilla-counter-width}
\end{equation}
\texttt{VanillaBpDecoder} loads one frame, drives
\texttt{core.io.bias=core.io.prior}, holds
\texttt{restartMessages=false}, and asserts \texttt{step.valid} throughout
\texttt{sRun}.  For completed iteration $t$, it drives the ramp control
\begin{equation}
 u_t=\min(t,2^{w_u}-1),
 \label{eq:scale-control}
\end{equation}
where $w_u$ is \texttt{scale.controlBits}.  On every \texttt{core.io.commit},
it evaluates
\begin{equation}
 \operatorname{stop}_{\mathrm{BP}}(t)=
 \operatorname{converged}(t)\lor[t=T_{\mathrm{frame}}].
 \label{eq:vanilla-stop}
\end{equation}
The terminal commit is copied to the common result; success and
\texttt{solutionsFound} are one exactly when it converged, and
\texttt{legsExecuted}=1.  If \texttt{runtimeIterationLimit=true}, a separate
\texttt{iterationLimit} port supplies $T_{\mathrm{frame}}$ for each accepted
frame subject to $1\leq T_{\mathrm{frame}}\leq T_{\max}$; otherwise the
elaborated capacity is used.  The implementation is in
\texttt{VanillaBpDecoder} in \path{BpDecoder.scala}; the independent executable
equation is \texttt{Reference.runVanilla} in \path{Reference.scala}.

\subsection{Relay coefficient generation}
\texttt{RelayCoefficientSource} in \path{RelayCoefficients.scala} defines a
replaceable source with \texttt{frameStart}, \texttt{advanceLeg}, and one
$b_{\ell,j}$ output per variable.  The production
\texttt{LfsrRelayCoefficients} uses one 16-bit Galois LFSR per variable.  Lane
$j$ is initialized on every frame to
\begin{equation}
 s_{0,j}=1+\left((o+32768j)\bmod 65535\right),
 \label{eq:lfsr-seed}
\end{equation}
where $o$ is \texttt{seedOffset}.  Leg zero uses $b_{0,j}=7$ without consuming
an LFSR word.  At each later-leg transition,
\begin{align}
 s_{\ell,j}&=(s_{\ell-1,j}\mathbin{\texttt{>>}}1)
 \oplus\begin{cases}
 \mathtt{0xB400},&s_{\ell-1,j}[0]=1,\\
 0,&s_{\ell-1,j}[0]=0,
 \end{cases}\label{eq:lfsr-step}\\
 b_{\ell,j}&=3+(s_{\ell,j}\mathbin{\&}7),
 \qquad \ell\geq1.\label{eq:lfsr-beta}
\end{align}
Thus later legs use $b\in[3,10]$, or effective
$\gamma=1-b/8\in[-0.25,0.625]$.  Coefficients remain constant inside a leg and
the complete sequence repeats across frames by design.  The independent Scala
implementation is \texttt{Reference.lfsrBeta}.

\subsection{Relay-BP controller}
\texttt{RelayBpConfig} in \path{RelayBpDecoder.scala} adds initial-leg budget
$T_0$, later-leg budget $T_r$, total hardware-leg capacity $L_{\max}$,
solution target $S$, seed offset, and Relay format.  Its maximum possible
iteration count is
\begin{equation}
 I_{\max}=T_0+(L_{\max}-1)T_r.
 \label{eq:relay-maximum-iterations}
\end{equation}
The derived counter widths are the positive bit lengths of $I_{\max}$,
$\max(T_0,T_r)$, $L_{\max}$, and $S$.  The solution-score width is the positive
bit length of
\begin{equation}
 Q_{\max}=n(2^{w_m}-1).
 \label{eq:score-width}
\end{equation}

\texttt{RelayBpDecoder} instantiates the same \texttt{StaticTannerCore} as
vanilla BP, one coefficient source, and Eq.~\eqref{eq:relay-bias} for every
variable.  A leg stops on its first syndrome-consistent commit or at its local
budget:
\begin{equation}
 \operatorname{stop}_{\ell}(t)=
 \operatorname{converged}_{\ell}(t)\lor
 \begin{cases}
 [t=T_0],&\ell=0,\\
 [t=T_r],&\ell>0.
 \end{cases}
 \label{eq:relay-leg-stop}
\end{equation}
Every converged leg increments the solution count, including a duplicate
solution.  A converged candidate is ranked by
\begin{equation}
 Q(\hat e)=\sum_{j=0}^{n-1}\Lambda_j\hat e_j.
 \label{eq:relay-score}
\end{equation}
The candidate replaces the stored best solution only when its score is
strictly smaller, so an equal-score tie retains the earliest solution.  The
decoder stops when the updated solution count reaches $S$ or when no configured
leg remains.  Otherwise, on the terminal VNU edge, \texttt{advanceLeg} both
advances the coefficients and asserts \texttt{restartMessages}.  The marginal
and decision from that commit are retained, V2C messages return to the immutable
priors, the local iteration counter becomes zero, and Eq.~\eqref{eq:ramp-alpha}
restarts at $t=1$ on the next CNU clock.  Hence a leg transition consumes zero
additional clocks.

On success, \texttt{correction} is the best converged candidate and the
reported residual is forced to zero because that candidate was checked when it
was recorded.  \texttt{marginal} remains the final executed leg's soft state,
which is the intended handoff to a later post-processor.  On total failure,
correction, marginal, and residual are one coherent snapshot of the last
commit.  Optional \texttt{runtimeLegLimit} adds a per-frame bound
$1\leq L_{\mathrm{frame}}\leq L_{\max}$ without re-elaboration.
\texttt{RelayLegTrace} emits the zero-based leg index, executed local
iterations, and convergence flag on every terminal commit; it is nonblocking.

\subsection{Data-structure delta from vanilla to Relay BP}
Table~\ref{tab:relay-delta} lists the additions needed to transform the vanilla
controller into Relay BP.  The Tanner graph, static connections, CNU/VNU
payloads, arithmetic leaf, and two-clock schedule are unchanged.

\begin{table*}[t]
\caption{Concrete additions from vanilla BP to Relay BP.}
\label{tab:relay-delta}
\centering
\scriptsize
\begin{tabular}{@{}p{0.18\textwidth}p{0.34\textwidth}p{0.40\textwidth}@{}}
\toprule
Concern & Vanilla representation & Relay addition and source symbol \\
\midrule
Configuration & \texttt{VanillaBpConfig}: nodes, $T_{\max}$, $q$, scale & \texttt{RelayBpConfig}: $T_0,T_r,L_{\max},S$, seed offset, \texttt{RelayFormat}, and derived local/leg/solution/score widths \\
Core inputs & \texttt{bias=prior}; restart false & Per-variable \texttt{RelayBias(...)} vector; \texttt{restartMessages=advanceLeg} \\
Controller state & \texttt{iteration}, \texttt{frameLimit}, terminal \texttt{outcome} & \texttt{leg}, \texttt{frameLegLimit}, \texttt{localIteration}, \texttt{totalIterations}, \texttt{solutions}, \texttt{haveBest}, \texttt{bestScore}, and \texttt{bestCorrection} \\
Coefficient state & None & Abstract \texttt{RelayCoefficientSource}; production per-VNU LFSR state and coefficient registers in \texttt{LfsrRelayCoefficients} \\
Result metadata & Common result with one-bit leg/solution counts & Same \texttt{BpDecoderResult} type parameterized with wider \texttt{legsExecuted} and \texttt{solutionsFound} fields \\
Observability & Frame result only & \texttt{RelayLegTrace} \texttt{Valid} event at each terminal VNU commit \\
Artifact ABI & Runtime iteration limit & Runtime total-leg limit plus leg/solution outputs and trace fields \\
\bottomrule
\end{tabular}
\end{table*}

\section{Legacy \texorpdfstring{$2\times2$}{2 x 2} Teaching Top}
\texttt{MinSumIteration2x2} predates the generic graph-generated top and remains
as the shortest registered wiring example for
\begin{equation}
 H=\begin{bmatrix}1&1\\1&1\end{bmatrix}.
\end{equation}
It accepts check-major input edges, instantiates two resetful
\texttt{CheckNodeUnit}s followed by two resetful \texttt{VariableNodeUnit}s,
and produces one iteration after two cycles.  It has no internal iterative
state or graph object.  New decoder work should normally start from
\texttt{StaticTannerDatapath}; the $2\times2$ top is retained for focused
pedagogy and regression coverage.

\section{Independent Systolic Gauss--Jordan Subsystem}
The \path{GaussJordan/} subtree intentionally imports no BP protocol or decoder
configuration type.  It demonstrates how a later post-processing stage can
remain separately elaborated, tested, emitted, and scheduled.

\texttt{GjOpcode} fixes the legacy two-bit encoding:
\texttt{Pass}=00, \texttt{Swap}=01, \texttt{Add}=10, and
\texttt{Lock}=11.  For stored bit $r$ and input $d$, \texttt{PeCol} passes $d$,
swaps or locks by emitting $r$ and storing $d$, and adds by emitting
$r\oplus d$.  \texttt{PeDiag} gives
\texttt{reduce\_sig\_i}$\land r$ priority and emits \texttt{Swap}; otherwise a
zero passes, the first one locks state, and subsequent ones emit
\texttt{Add}.  State, data, and opcode registers have active-high synchronous
reset and hold on disabled cycles.  The diagonal reduction output is
combinational, reset-masked, and deliberately not enable-gated.

\texttt{TrapezoidMeshConfig} validates positive $n$, lifted-column count, and
reduction-hop delay and derives the total column count.  The elaborated mesh
ties off its lower-left triangle, moves data vertically, moves opcodes
horizontally, and delays reduction only between diagonal cells on enabled
cycles.  State exports pack $A$ and lifted $B$ in row-major order with
coordinate zero at the least-significant bit.  \texttt{FlatIO} and explicit
\texttt{desiredName}s preserve the legacy external ABI; the historical top
spelling \texttt{trapeziod\_mesh} is intentionally retained.  A streamed row
feeder, row count, and decoder/post-processor controller are not present.

\section{Emission and Artifact Organization}
\subsection{Build contract}
The repository pins Mill 1.1.2, Scala 2.13.18, Chisel 7.7.0, and ScalaTest
3.2.19 in \path{build.mill}.  The matching Chisel compiler plugin is enabled,
and warning-oriented Scala compiler options include deprecation, feature, and
initialization checks.  The included Mill wrapper is the sole supported build
interface.  \path{.mill-jvm-opts} supplies the project root required by the
Chisel test-directory flow.  Continuous integration installs Java 17 and
Verilator and runs the complete Mill test target for main-branch or tag pushes,
pull requests, and manual dispatches.

\subsection{Named generator tops}
\texttt{Generate.main} maps a short research target name to a zero-argument
module generator.  Each invocation writes one FIRRTL-dialect MLIR file and one
SystemVerilog file under a requested directory.  The two artifacts come from
fresh but equivalently configured elaborations of the selected closure.  CIRCT
emission disables random initialization, strips debug information, and enables
default layer specialization so generated artifacts are deterministic and
compact.

\begin{table}[t]
\caption{Named emission targets.}
\label{tab:targets}
\centering
\scriptsize
\begin{tabular}{@{}p{0.33\columnwidth}p{0.60\columnwidth}@{}}
\toprule
Target & Emitted top \\
\midrule
\texttt{two-min} & \texttt{TwoMin} \\
\texttt{check-valls} & \texttt{CheckNode}, constant scaling \\
\texttt{check-relay} & \texttt{CheckNode}, ramp scaling \\
\texttt{variable} & \texttt{VariableNode} \\
\texttt{iteration} & legacy \texttt{MinSumIteration2x2} \\
\texttt{convergence} & Steane convergence checker \\
\texttt{static-steane} & full static Steane datapath \\
\texttt{vanilla-steane} & autonomous four-iteration vanilla BP \\
\texttt{relay-steane} & autonomous four-leg Relay BP \\
\texttt{bp-filtered-}\newline\texttt{osd0-steane} & composed BP and filtered OSD-0 \\
\texttt{neighbour-sorter} & standalone reliability sorter \\
\texttt{pe-col}, \texttt{pe-diag} & Gauss--Jordan leaf PEs \\
\texttt{trapezoid-mesh} & three-row example mesh \\
\bottomrule
\end{tabular}
\end{table}

\path{scripts/emit.sh} keeps emission separate from tests and accepts a target
and optional output directory.  \path{scripts/verify-rtl.sh} consumes an already
emitted SystemVerilog file and top name and performs standalone Verilator lint;
it does not elaborate Scala or prove behavior.  It disables only warning
classes caused by CIRCT's single-file naming, unused portions of intentional
fixed-point intermediates, and intentionally open boundary pins.  Generated
directories are excluded from version control.  A result-bearing emitted
artifact should instead be archived immutably with graph/configuration input,
Git revision, and Scala, Chisel, CIRCT, Verilator, synthesis, and FPGA tool
versions.

\section{Verification Methodology}
The verification ladder crosses semantic boundaries deliberately.  It does not
use emitted RTL text as the golden specification.

\subsection{Pure Scala arithmetic oracle}
\path{Reference.scala} defines test-only messages and decoder state using
ordinary \texttt{Int} and \texttt{Boolean}.  Its two-minimum oracle sorts, its
check and variable functions use direct integer equations, and its graph
iteration follows the canonical node plan.  It calls no Chisel arithmetic
helper.  It does share validated configuration and graph metadata with the DUT;
separate graph tests therefore lock topology and port mapping before the model
is used for graph-level comparisons.

\path{Harnesses.scala} wraps combinational \texttt{RawModule}s in small clocked
test tops where required by ChiselSim.  These harnesses add no alternative
algorithm; they expose leaf IO to the simulator.

\subsection{Executable test inventory}
Table~\ref{tab:tests} maps each maintained specification to its claim.  This is
also the recommended order for debugging a failed graph-level test.

\begin{table*}[t]
\caption{Verification inventory and implemented claim.}
\label{tab:tests}
\centering
\scriptsize
\begin{tabular}{@{}p{0.235\textwidth}p{0.69\textwidth}@{}}
\toprule
Specification or executable & Claim established \\
\midrule
\path{ConfigSpec.scala} & Relay defaults are locked; representative LLR probabilities quantize and saturate as specified. \\
\path{TannerGraphSpec.scala} & Rows normalize, column/check/variable views derive correctly, edge IDs are stable, and malformed adjacency is rejected. \\
\path{TannerNodeGraphsSpec.scala} & Exact Steane check and variable views, degrees, global-edge bijection, and all local port bindings. \\
\path{TwoMinSpec.scala} & Exhaustive $8^4$ comparison for four three-bit inputs, including all duplicate-minimum cases. \\
\path{CheckNodeSpec.scala} & Exhaustive unscaled degree-three behavior over all two-bit magnitudes, signs, and syndromes; directed Valls and ramp scaling. \\
\path{VariableNodeSpec.scala} & Fixed-seed randomized base VNU comparison, full eight-bit hard-decision range, saturation, and registered feedback. \\
\path{RelayBiasSpec.scala} & Reduced partial-product truncation, signs, saturation, and five-bit paper-profile bias against an independent integer model. \\
\path{RelayCoefficientsSpec.scala} & Lane seeds, Galois recurrence, initial coefficient, leg stability, advancement, and repeatable frame reseeding. \\
\path{ConvergenceCheckerSpec.scala} & All $2^7$ Steane estimates crossed with all $2^3$ syndromes for $H\hat e\oplus S$ and convergence. \\
\path{IterationSpec.scala} & The legacy $2\times2$ result and its two-cycle validity protocol. \\
\path{StaticTannerDatapathSpec.scala} & Step-before-load rejection, prior-to-edge initialization, all eight Steane syndromes, four consecutive iterations, two-cycle timing, marginals, decisions, residual, and convergence against the Scala oracle. \\
\path{BpDecodersSpec.scala} & Vanilla and Relay frame control, runtime limits, output stalls, persistent marginals, message restart, zero-cycle leg transitions, scoring/ties, failure snapshot, production LFSRs, and external oracle fixture. \\
\path{SparseBitmaskStreamerSpec.scala} & Ascending sparse-one output, empty masks, banking, and backpressure stability. \\
\path{BpOnlyArtifactSpec.scala} & BP-only wrapper, independent model, empty/sparse corrections, exact always-ready cycle equation, and output stalls. \\
\path{RelayBpArtifactSpec} in \path{RelayBpExperiment.scala} & Default Relay-5 profile, independent Relay result, per-leg trace, sparse stream, final soft state, and exact cycle equations. \\
\path{EndToEndSpec.scala} & Strict file-schema, graph-dimension, prior-range, syndrome, and iteration validation before elaboration. \\
\path{GaussJordanPESpec.scala} & Opcode encodings, all PE transitions, synchronous reset, disabled hold, and combinational reduction forwarding. \\
\path{TrapezoidMeshSpec.scala} & Configuration rejection, boundary behavior, an independent cycle oracle under deterministic and random traffic, row-major packing, hold/reset, and reduction delays one through three. \\
\bottomrule
\end{tabular}
\end{table*}

The complete ScalaTest suite is run by \path{scripts/test.sh} and by continuous
integration.  ChiselSim exercises the generated hardware with Verilator.  The
standalone lint script is a separate, weaker check: it establishes that a named
emitted artifact is accepted by Verilator with the selected warnings, not that
its values match the algorithm.

\section{Shared BB144 Experimental Contract}
\subsection{Code, circuit, detector model, and samples}
The two maintained decoder experiments deliberately share one problem source:
\path{benchmarks/bb144/bb144.py}.  It reconstructs the pinned gross-code
$[[144,12,12]]$ CSS matrices from $12\times6$ bivariate circulants.  If
$X=P_{12}\otimes I_6$ and $Y=I_{12}\otimes P_6$, the implemented polynomials
and parity checks are
\begin{align}
 A&=X^3+Y+Y^2,&B&=Y^3+X+X^2,\label{eq:bb-polynomials}\\
 H_X&=[A\mid B],&H_Z&=[B^{\mathsf T}\mid A^{\mathsf T}],
 \label{eq:bb-checks}
\end{align}
with addition over $\mathrm{GF}(2)$.  \texttt{code\_data} checks the resulting
$72\times144$ matrices, CSS commutation, 12 logical rows, and the upstream
commit recorded by \texttt{UPSTREAM\_COMMIT}.

\texttt{circuit} constructs the 12-noisy-cycle memory circuit plus two
zero-noise closure rounds.  At physical error probability $p$, reset and
measurement flips occur with probability $p$, one-qubit noisy locations use a
uniform nonidentity Pauli conditioned on an event of probability $p$, and
two-qubit noisy locations analogously use one of 15 nonidentity Paulis.
\texttt{sample\_physical} independently implements this circuit-level Pauli
process in vectorized NumPy and returns the Z-check detector differences and 12
actual logical-observable bits.  Probability point $a$ uses deterministic seed
\begin{equation}
 z_a=z_0+a,\qquad z_0=14412
 \label{eq:sample-seed}
\end{equation}
unless the command line overrides $z_0$.

\texttt{z\_check\_model} also obtains an undecomposed Stim detector error model
(DEM).  A decoder column is keyed by its complete detector support and logical
support.  Independent effects with the same key are merged by the XOR law
\begin{equation}
 p_{a\oplus b}=p_a(1-p_b)+p_b(1-p_a).
 \label{eq:xor-probability}
\end{equation}
Equal detector supports with different logical supports remain distinct.
For DEM-column probability $p_j\in(0,1/2)$, \texttt{quantize} computes
\begin{equation}
 \Lambda_j=\min\!\left(15,
 \left\lfloor2\ln\frac{1-p_j}{p_j}+\frac12\right\rfloor\right).
 \label{eq:prior-quantization}
\end{equation}

The min-sum CNU requires degree at least two, so
\texttt{reduce\_degree\_one} solves every singleton row exactly.  If singleton
check $i$ fixes column $j$ then $e_j=S_i$.  Let $F$ be the fixed-column set and
$K$ the retained checks.  \texttt{reduce\_sample} forms
\begin{align}
 S'_i&=S_i\oplus\bigoplus_{j\in\mathcal N(i)\cap F}e_j,
 &&i\in K,\label{eq:reduced-syndrome}\\
 o'_q&=o_q\oplus\bigoplus_{j\in\mathcal L(q)\cap F}e_j,
 &&0\leq q<12.\label{eq:reduced-logical}
\end{align}
It then deletes the fixed variables and singleton checks and remaps column
indices.  The current code asserts the reduced dimensions
\begin{equation}
 (m,n,E,k_L)=(936,8784,30672,12).
 \label{eq:bb144-reduced-dimensions}
\end{equation}
For every requested $p$, topology must match the first point; only priors and
sampled frames vary.  \texttt{prepare} archives the circuit, one DEM per point,
\texttt{sweep.json}, the complete \texttt{benchmark.txt} stream, and the first
shot at each point as \texttt{verify/p*/problem.json} plus
\texttt{actual.json}.  This separation makes the static graph and runtime data
auditable.

\subsection{Artifact output and exact cycle contract}
\texttt{BpOnlyArtifact} in \path{BpOnlyExperiment.scala} and
\texttt{RelayBpArtifact} in \path{RelayBpExperiment.scala} expose the same
packed benchmark ABI.  Check zero and variable zero occupy the least-significant
fields.  Both wrappers include unused scope and OSD-compatible accounting ports
so the common C++ transport can also serve composed decoders; for the two
experiments here, \texttt{osdCycles}, \texttt{selected}, \texttt{activeRows},
and \texttt{solverCycles} are identically zero.

The dense decoder correction is snapshotted by
\texttt{SparseBitmaskStreamer} in \path{SparseBitmaskStreamer.scala}.  With bank
width $W_b=64$, correction support
$\mathcal C=\{j:\hat e_j=1\}$, correction weight $w=|\mathcal C|$, and occupied
bank count
\begin{equation}
 b=\left|\left\{\left\lfloor\frac{j}{W_b}\right\rfloor:j\in\mathcal C\right\}\right|,
 \label{eq:occupied-banks}
\end{equation}
the stream emits the members of $\mathcal C$ in strictly increasing order.
An empty correction emits no placeholder; \texttt{resultValid} completes the
frame.  With an always-ready consumer, the measured wrapper latency from the
accepted input edge through result visibility is
\begin{equation}
 C_{\mathrm{frame}}=2I+1+b+w.
 \label{eq:sparse-frame-clocks}
\end{equation}
The $2I$ term is Eq.~\eqref{eq:bp-clocks}, one clock transfers the decoder result
into the streamer, each occupied bank costs one selection clock, and each set
bit costs one transfer clock.  If correction output is stalled for $d$ clocks,
$C_{\mathrm{frame}}$ increases to $2I+1+b+w+d$.  Stalling
\texttt{resultReady} after \texttt{resultValid} holds the completed frame and
does not increment the reported counter.  \texttt{checkSparseTiming} in
\path{sim_main.cpp}, \texttt{BpOnlyArtifactSpec}, and
\texttt{RelayBpArtifactSpec} enforce these equations.

\section{Experiment 1: Vanilla BP on BB144}
\subsection{Decoder profile}
The runnable entry point is
\path{examples/BivariateBicycle144BpOnly/run.sh}.  It elaborates
\texttt{VanillaBpDecoder} through \texttt{BpOnlyArtifact} with
\begin{equation}
 (w_m,w_a,T_{\max})=(4,7,30),
 \qquad \alpha_t\text{ given by Eq.~\eqref{eq:ramp-alpha}}.
 \label{eq:bp-only-profile}
\end{equation}
The decoder stops at the first converged commit or after 30 iterations.  A
nonconverged terminal hard decision is still streamed for diagnosis, but status
is numeric 2 (\texttt{bp\_nonconverged}) and the experiment classifies the shot
as a failure.  \texttt{BpOnlyArtifact} enables the optional runtime
\texttt{iterationLimit} port; an emitted 30-iteration-capacity artifact can
therefore execute any per-frame prefix from 1 through 30 without re-elaboration.
The standalone BP-only sweep drives 30 unless the budget meta-experiment is
selected.

The script defaults to 1000 deterministic shots at each
$p\in\{0.001,0.002,\ldots,0.007\}$, while \texttt{P\_VALUES} and the second
positional argument can override those choices.  These defaults define a
maintained executable, not a statistical sufficiency claim.

\subsection{Traceable workflow}
The BP-only workflow is the following ordered dataflow; every arrow names a
materialized input or output rather than an implicit in-memory handoff.
\begin{enumerate}
  \item \texttt{bb144.py prepare} writes \path{source/benchmark.txt},
        \path{source/sweep.json}, circuit/DEM provenance, and one
        \path{source/verify/p*/problem.json} per probability point.
  \item \texttt{BpOnlyExperiment.main} parses one problem with
        \texttt{DecoderProblem.read}, writes \path{artifact/config.txt} and
        versioned \path{artifact/config.json}, evaluates
        \texttt{BpOnlyExperiment.golden}, and emits the specialized
        \path{artifact/rtl/BpOnlyArtifact.sv} through
        \texttt{Generate.emitSystemVerilogFiles}.
  \item \texttt{BpOnlySweepGolden.main} runs the independent
        \texttt{Reference.iterate} loop for the first shot at every $p$ and
        writes \path{verification/p*/golden.txt}.  It stops on exactly the same
        predicate as Eq.~\eqref{eq:vanilla-stop} but shares no Chisel arithmetic.
  \item Verilator compiles the saved SystemVerilog together with
        \path{benchmarks/bb144/sim_main.cpp} under
        \texttt{BP\_ONLY\_ARTIFACT}.  Exact mode runs every golden before any
        performance sweep is accepted.
  \item \texttt{adapters/chipsldpc.py} splits the immutable benchmark stream
        into shot shards and records SHA-256 hashes of each input, the simulator,
        and \path{artifact/config.txt} in \path{parallel/jobs.json}.
  \item \texttt{workflow.py run} executes independent Verilator processes;
        \texttt{workflow.py merge} revalidates hashes, row counts, probability
        labels, and unique shot indices before writing \path{raw/shots.csv}.
  \item \texttt{report.py} validates timing sums and corpus completeness, then
        writes \path{results/results.csv}, \path{results/timing.csv},
        \path{results/results.json}, and the logical-failure/cycle plots.
\end{enumerate}
The emitted RTL is generated once because the graph is common across $p$;
priors and syndromes are runtime inputs.  Comparing
\path{source/benchmark.txt} across decoder runs proves whether the two
experiments consumed an identical shot corpus.

\section{Experiment 2: Relay-BP-S on BB144}
\subsection{Relay-5 parameter profile}
The runnable entry point is
\path{examples/BivariateBicycle144RelayBp/run.sh}.  Its default experiment
profile is
\begin{equation}
 (w_m,w_a,f,T_0,T_r,R,S,o)=(4,5,3,80,60,600,5,0),
 \label{eq:relay-experiment-profile}
\end{equation}
where $R$ is the number of noninitial legs and $S$ is the
target number of converged legs.  \texttt{RelayBpExperimentProfile.decoderConfig}
maps this convention to the lower-level hardware field
\begin{equation}
 L_{\max}=R+1=601,
 \qquad I_{\max}=T_0+RT_r=36080.
 \label{eq:relay-profile-capacity}
\end{equation}
This resolves the source paper's total-leg wording in favor of its Algorithm~1
and reference-software loop convention~\cite{maurer2025}.  Setting
\texttt{RELAY\_R=599} instead produces exactly 600 total legs.

Leg zero uses $b=7$, hence $\gamma_0=1-7/8=0.125$.  Later coefficients from
Eq.~\eqref{eq:lfsr-beta} give the effective quantized interval
$\gamma\in[-0.25,0.625]$ for the requested paper interval
$[-0.24,0.66]$.  The LFSR polynomial/seed mapping is a reproducible
\texttt{chipsLDPC} choice because the paper does not publish those details; it
is serialized in \path{artifact/config.json}.  Every frame reseeds, so all
shots see the same coefficient schedule and differ only in their priors and
syndromes.

The script defaults to 10,000 shots at
$p\in\{0.001,0.002,0.003,0.004\}$; the example README shows a smaller explicit
1000-shot invocation.  \texttt{P\_VALUES}, positional shot count, and
\texttt{RELAY\_T0}, \texttt{RELAY\_TR}, \texttt{RELAY\_R},
\texttt{RELAY\_S}, and \texttt{RELAY\_SEED\_OFFSET} parameterize a run.  This
is the Z-check detector sector of the 12-cycle memory experiment.  A full XZ
curve would require running and combining the independently decoded X-check
sector; the current result must not be labeled as that full-code curve.

\subsection{Workflow and Relay-specific records}
The Relay script follows the same seven materialized stages as Experiment~1,
with the following substitutions:
\begin{itemize}
  \item \texttt{RelayBpExperiment.main} emits \texttt{RelayBpArtifact} and a
        versioned Relay configuration; \texttt{RelayBpSweepGolden.main} calls
        \texttt{Reference.runRelay} and \texttt{Reference.lfsrBeta};
  \item Verilator compiles \path{sim_main.cpp} under
        \texttt{RELAY\_BP\_ARTIFACT};
  \item the runtime \texttt{legLimit} selects a prefix of the emitted 601-leg
        capacity; and
  \item each terminal VNU commit adds the tuple
        $(\ell,I_\ell,C_\ell,c_\ell)$ to the trace, where
\begin{align}
 C_\ell&=2I_\ell,\\
 I&=\sum_{\ell=0}^{L-1}I_\ell,\\
 C_{\mathrm{BP}}&=\sum_{\ell=0}^{L-1}C_\ell=2I,\\
 N_{\mathrm{sol}}&=\sum_{\ell=0}^{L-1}c_\ell.
 \label{eq:relay-trace-invariants}
\end{align}
\end{itemize}
Here $c_\ell\in\{0,1\}$ records convergence, not logical correctness.  The
trace is written compactly into each raw CSV row and expanded by
\texttt{report.py} into \path{results/relay_leg_timing.csv},
\path{relay_leg_timing_summary.csv}, and \path{relay_leg_timing.png}.

If \texttt{RELAY\_REF\_DIR} names the pinned editable \texttt{trmue/relay}
checkout, the script first runs \path{benchmarks/oracles/trmue_relay.py} on the
pinned repetition-code coefficient fixture and archives
\path{verification/trmue-relay.json}.  This optional cross-implementation
check is separate from, and does not replace, the mandatory Scala golden and
exact emitted-RTL checks.

\section{Exact RTL Simulation and Statistical Reduction}
\subsection{The dual-mode \texttt{sim\_main.cpp} harness}
\path{benchmarks/bb144/sim_main.cpp} is the behavioral boundary between the
saved SystemVerilog artifact and all accepted experiment data.  A compile-time
macro selects \texttt{VBpOnlyArtifact} or \texttt{VRelayBpArtifact}; the
arithmetic and controller remain entirely in emitted RTL.  Template helpers
\texttt{setBit}, \texttt{getBit}, \texttt{setField}, and
\texttt{signedField} handle both scalar and Verilator wide ports and decode
two's-complement soft fields without assuming a C++ layout for wide signals.

\texttt{readConfig} parses \path{artifact/config.txt} and rejects mismatched
ABI versions, dimensions, widths, invalid Relay counts, or excessive capacity.
For vanilla BP, a requested iteration budget $B$ maps directly to the hardware
runtime limit.  For Relay BP it must satisfy
\begin{equation}
 B=T_0+rT_r,\qquad r\in\mathbb Z_{\geq0},
 \label{eq:relay-runtime-budget}
\end{equation}
and \texttt{decoderRuntimeLimit} drives
\begin{equation}
 L_{\mathrm{frame}}=1+\frac{B-T_0}{T_r}.
 \label{eq:budget-to-legs}
\end{equation}
Both mappings are checked against the capacity recorded beside the RTL.

\texttt{decode} performs one complete hardware transaction in five explicit
phases:
\begin{enumerate}
  \item clear all packed inputs, set the runtime limit, and pack syndrome bit
        $i$ at offset $i$ and prior $j$ at offset $jw_m$;
  \item assert \texttt{inputValid}, require \texttt{inputReady}, clock the
        accepting edge, and then deassert \texttt{inputValid};
  \item on each subsequent clock, collect every sparse correction transfer and
        Relay leg-trace event until \texttt{resultValid};
  \item reject any repeated or out-of-range correction index, require
        \texttt{correctionLast} only on the final nonempty transfer, reconstruct
        the dense correction, and sign-extend each packed $w_a$-bit soft output;
  \item compare the RTL cycle counter with the harness's observed clock count,
        accept the result, and clock once more to return the wrapper to idle.
\end{enumerate}
For a requested iteration bound $B$ and configured OSD prefixes $p_k$, the
benchmark-mode deadlock watchdog is
\begin{equation}
 C_{\max}=4n+2m+3\sum_k p_k+2B+128.
 \label{eq:sim-watchdog}
\end{equation}
Exact mode substitutes the golden iteration count for $B$.  This bound is only
a failure detector; it is never reported as decoder latency.

\texttt{checkSparseTiming} enforces Eqs.~\eqref{eq:bp-clocks} and
\eqref{eq:sparse-frame-clocks}.  For these BP-only artifacts,
\texttt{breakdown} consequently records
\begin{equation}
 C_{\mathrm{dispatch}}=1,\qquad
 C_{\mathrm{BP\ output}}=b+w,
 \label{eq:sim-breakdown}
\end{equation}
with all sort, solver, and OSD terms zero.  \texttt{checkRelayTiming} enforces ordered leg
indices, per-leg limits, Eq.~\eqref{eq:relay-trace-invariants}, and consistency
of the solution count.  A trace or timing mismatch is therefore a failed run,
not a post-processing warning.

The executable has two modes:
\begin{itemize}
  \item \texttt{exactMain(golden,result)} reads an independent Scala golden,
        calls \texttt{decode} at the artifact's full configured capacity, and
        compares status, iteration and cycle counts, soft vector, sparse stream,
        and all Relay metadata.  Only after every field matches does it write
        \texttt{result.json}.
  \item \texttt{benchmarkMain(input,csv,config,budget)} reads the BB144 stream,
        runs every shot on the same already-verified RTL object, checks all
        transport/timing invariants, evaluates logical observables, and writes
        one raw CSV row per shot.  The optional budget selects the runtime path
        in Eqs.~\eqref{eq:relay-runtime-budget}--\eqref{eq:budget-to-legs}.
\end{itemize}
The Scala golden files are produced by
\texttt{BpOnlyExperiment.golden} or \texttt{RelayBpExperiment.golden}.  Their
configuration line, syndrome, priors, terminal status/counters, soft state, and
sparse correction are a compact verification protocol; the Relay record also
contains every expected leg tuple.

\subsection{Logical classification and reported estimators}
For logical-support row $\mathcal L(q)$, \texttt{benchmarkMain} predicts
\begin{equation}
 \hat o_q=\bigoplus_{j\in\mathcal L(q)}\hat e_j.
 \label{eq:predicted-observable}
\end{equation}
Let $a_s$ be the decoder-success predicate for shot $s$ (vanilla convergence or
at least one Relay solution), and let $o_s$ be the 12-bit physical-sample
observable vector after Eq.~\eqref{eq:reduced-logical}.  The recorded word
failure indicator is
\begin{equation}
 F_s=[\neg a_s]\lor[\hat o_s\neq o_s].
 \label{eq:word-failure-event}
\end{equation}
For $N$ shots at one $p$, \texttt{report.py} computes
\begin{align}
 \widehat P_{\mathrm{word}}&=\frac1N\sum_{s=1}^{N}F_s,
 \label{eq:word-failure-estimator}\\
 \widehat p_{L/\mathrm{cycle}}&=
 1-(1-\widehat P_{\mathrm{word}})^{1/12},
 \label{eq:per-cycle-rate}\\
 \overline C&=\frac1N\sum_{s=1}^{N}C_s.
 \label{eq:mean-cycles}
\end{align}
Equation~\eqref{eq:per-cycle-rate} converts a 12-QEC-cycle experiment word
probability under a constant independent-rate interpretation; it does not
change the underlying shot-level failure definition.  The conditional BP plot
uses only shots with successful BP status in both numerator and denominator and
counts $[\hat o_s\neq o_s]$ among that subset.

For $x=\sum_sF_s$, $\hat p=x/N$, and
$z=1.959963984540054$, the reported Wilson 95\% interval is
\begin{align}
 d&=1+\frac{z^2}{N},\\
 c&=\frac{\hat p+z^2/(2N)}{d},\\
 r&=\frac{z}{d}\sqrt{\frac{x(N-x)}{N^3}+\frac{z^2}{4N^2}},\\
 [p_-,p_+]&=[\max(0,c-r),\min(1,c+r)].
 \label{eq:wilson}
\end{align}
Cycle medians, means, maxima, and nearest-rank 95th percentiles are computed
from measured $C_s$.  Stage totals are accepted only if their sum equals every
shot's \texttt{total\_cycles}.  Relay summaries additionally aggregate the
validated leg traces by probability and one-based presentation leg number.

\subsection{Hashed parallel execution}
\path{benchmarks/rtl_shards/adapters/chipsldpc.py} writes one input file per
$(p,\text{shot range})$ and a manifest containing SHA-256 digests of every
input, the simulator executable, and its configuration.  With a budget sweep,
the same input shard is referenced by one job per budget; no resampling occurs.
\texttt{workflow.py run} skips only an already present output whose source hash,
row count, local indices, and group labels all validate.  It writes through a
temporary file and atomically renames on success.  \texttt{workflow.py merge}
repeats artifact and row validation, converts local to global shot indices, and
rejects duplicates before producing the raw corpus.  Host wall time and
throughput are recorded separately from architecture clocks.

\subsection{Runtime iteration-budget meta-experiment}
The iteration-budget example's \texttt{run.sh} is an analysis
workflow over either Experiment~1 or Experiment~2, not a third decoder.  It
prepares one probability point and one deterministic shot corpus, emits one
artifact at the maximum requested capacity, and runs every shot at every budget
through the runtime-limit port.  Vanilla accepts any positive budget not
exceeding capacity.  Relay accepts only Eq.~\eqref{eq:relay-runtime-budget}; the
script converts the largest budget into the emitted $R$ capacity and the C++
harness converts each individual budget back into a total-leg limit.

For budget $B$ and executed iterations $I_{B,s}$,
\path{benchmarks/bb144/budget_report.py} reports
\begin{align}
 \overline I_B&=\frac1N\sum_{s=1}^{N}I_{B,s},\\
 h_B&=1.96\frac{s_B}{\sqrt N},\\
 I_{B,95}&=I_{B,(\lceil0.95N\rceil)},\\
 \widehat P_{\mathrm{word},B}&=\frac1N\sum_{s=1}^{N}F_{B,s},
 \label{eq:budget-metrics}
\end{align}
where $s_B$ is the sample standard deviation and $I_{B,(k)}$ is the $k$th
ordered value.  It rejects an incomplete budget group, duplicate shot index, or
any $I_{B,s}>B$.  The plotted point is
$(\overline I_B,\widehat P_{\mathrm{word},B})$ with the mean-iteration normal
half-width and Wilson failure interval; zero observed failures are plotted at
their Wilson upper bound.  Outputs are
\path{results/budget_sweep.csv}, \path{budget_sweep.json}, and
\path{logical_failure_vs_average_iterations.{png,pdf}}.

\subsection{What simulation does not establish}
Behavioral tests establish finite sets of values, protocol conditions, and
architecture-clock latency.  A completed many-shot run estimates failure under
its recorded sample distribution and reports finite-sample intervals; it does
not prove the population failure rate.  Neither test class establishes maximum
clock frequency, critical path, lookup-table or register count, routing
congestion, memory use, power, or throughput across independently resident
shots.  Converting architecture clocks into time requires a named FPGA,
synthesis/place-and-route version, constraints, and achieved clock period.

\section{Repository Map and Reproduction}
The maintained repository is deliberately shallow:
\begin{itemize}
  \item \path{src/main/scala/chipsldpc/} contains pure elaboration records,
        hardware generators, and the named emission entry point;
  \item \path{src/test/scala/chipsldpc/} contains independent arithmetic models,
        simulation harnesses, specifications, and golden-record generation;
  \item \path{scripts/} contains one concise action per script: test, emit,
        lint an emitted RTL artifact, or create the BB144 Python environment;
  \item \path{examples/} contains the maintained BB144 BP-only and Relay-BP
        flows plus the runtime-budget meta-experiment;
  \item \path{benchmarks/bb144/} contains their shared circuit, sampling,
        Verilator transport, and reporting support;
  \item \path{build.mill}, the Mill wrapper, and \path{.mill-jvm-opts} define the
        reproducible build entry; and
  \item \path{paper/main.tex} is this evolving implementation record.
\end{itemize}
IDE metadata is not part of the architecture.  Generated \path{build/}, test,
and TeX auxiliary products are ignored and must not be cited as immutable
research outputs.

From the repository root, use one shell session for the commands below.  Run
the complete behavioral suite:
\begin{verbatim}
./scripts/test.sh
\end{verbatim}
Emit and lint the static decoder as two separate steps:
\begin{verbatim}
./scripts/emit.sh static-steane
./scripts/verify-rtl.sh \
  build/generated/static-steane/\
StaticTannerDatapath.sv \
  StaticTannerDatapath
\end{verbatim}
Optional leaf RTL inspection uses named emission targets rather than
duplicating their focused tests as examples:
\begin{verbatim}
./scripts/emit.sh pe-col
./scripts/emit.sh pe-diag
./scripts/emit.sh trapezoid-mesh
\end{verbatim}
Run the maintained BB144 vanilla-BP benchmark with:
\begin{verbatim}
./scripts/setup-bb144.sh
cd examples/BivariateBicycle144BpOnly
OUT=../../build/generated/bb144
./run.sh ${OUT}-bp-only 1000
\end{verbatim}
Run the Relay-BP experiment with:
\begin{verbatim}
cd ../BivariateBicycle144RelayBp
./run.sh ${OUT}-relay-bp 1000
\end{verbatim}
The default Relay parameters are Eq.~\eqref{eq:relay-experiment-profile}; the
environment overrides are listed in that example's README.  Run the budget
meta-experiment, selecting either \texttt{DECODER=vanilla} or
\texttt{DECODER=relay}, with:
\begin{verbatim}
cd ..
cd BivariateBicycle144IterationBudgetSweep
export DECODER=relay P_VALUE=0.003
./run.sh ${OUT}-relay-budget 10000
\end{verbatim}
Focused structure checks use Mill's \texttt{testOnly} target with the fully
qualified specification name.

When inspecting a new decoder change, the intended research loop is: change
one pure configuration or leaf rule; extend the independent Scala oracle and
focused specification; run that test; emit the smallest relevant top; inspect
the FIRRTL-dialect MLIR for widths and structure; inspect and lint the emitted
SystemVerilog; then run the exact artifact comparison if the static top is
affected.  This sequence localizes semantic, elaboration, lowering, and
artifact-integration failures without creating parallel implementations in
production source.

\section{Planned Extensions}
The next additions should preserve the same boundaries:
\begin{enumerate}
  \item synthesize and place the named BB144 artifacts, archive constraints and
        reports, and replace paper-clock conversions with achieved frequencies;
  \item bank or replicate state to evaluate throughput for multiple resident
        syndrome streams rather than sequential single-frame latency alone;
  \item implement and combine the complementary X-check sector under the same
        sample provenance before comparing against a full-code XZ curve;
  \item compare alternative, explicitly versioned coefficient generators and,
        if available, reproduce the external paper's unpublished PRNG mapping;
  \item connect the autonomous BP result to the existing filtered OSD-0 or
        Gauss--Jordan subsystem under one stable artifact ABI; and
  \item add immutable result manifests containing Git revision and exact Scala,
        Chisel, CIRCT, Verilator, Python, Stim, synthesis, and FPGA tool versions.
\end{enumerate}
A new policy should first be a small immutable Scala choice or a narrow bundle
extension.  Graph topology must remain elaboration-only, and benchmarking must
consume emitted artifacts rather than import the Chisel generator.

\section{Threats to Validity}
Small exhaustive configurations expose ties and fixed-width corner cases but do
not prove every degree/width combination.  The Scala oracle and DUT share the
validated graph metadata, although graph construction and port mapping have
separate structural tests; the comparison is strong regression evidence, not a
formal equivalence proof.  The BB144 study covers one reduced Z-check detector
sector and one circuit/noise construction.  It is not a full XZ gross-code
result, and exact degree-one preprocessing changes the hardware graph while
preserving the represented equations.  The deterministic coefficient schedule
supports reproducibility but is not evidence that the unpublished PRNG used by
the external FPGA work is identical.  Quantized-prior scoring is a hardware
selection rule, not proof of minimum physical correction weight.  Reusing the
same shots across decoder/budget variants improves paired comparison but does
not remove finite-sample uncertainty.  Finally, Verilator clock counts are
architectural; without place-and-route they imply neither frequency nor wall
clock latency on an FPGA.

\section{Conclusion}
The present \texttt{chipsLDPC} revision implements two complete autonomous
decoders over one statically wired, two-clock min-sum core.  Vanilla BP adds
only frame and iteration control.  Relay BP extends that controller with a
reduced memory-bias calculation, deterministic per-variable coefficients,
leg/local/total counters, persistent marginal state, message restart, candidate
scoring, and per-leg trace metadata; it does not duplicate the Tanner datapath.
Both variants are emitted as exact BB144 SystemVerilog artifacts, checked
against independent Scala outcomes, exercised on the same deterministic
circuit-level samples, and reduced through hash-validated raw records.  The
explicit equations and named source symbols in this manuscript separate
algorithmic behavior, architecture-clock accounting, statistical estimation,
and as-yet-unmeasured FPGA implementation claims.

\begin{thebibliography}{9}
\bibitem{valls2021}
J.~Valls, F.~Garcia-Herrero, N.~Raveendran, and B.~Vasi\'{c},
``Syndrome-based min-sum versus OSD-0 decoders: FPGA implementation and
analysis for quantum LDPC codes,'' \emph{IEEE Access}, vol.~9, 2021,
doi: 10.1109/ACCESS.2021.3118544.

\bibitem{maurer2025}
T.~Maurer et al., ``Real-time decoding of the gross code memory with FPGAs,''
arXiv:2510.21600, 2025.

\bibitem{steane1996}
A.~M. Steane, ``Error correcting codes in quantum theory,''
\emph{Phys. Rev. Lett.}, vol.~77, pp.~793--797, 1996,
doi: 10.1103/PhysRevLett.77.793.

\bibitem{chiselbook}
M.~Schoeberl, \emph{Digital Design with Chisel}, 6th ed., 2025.

\bibitem{chisel}
CHIPS Alliance, ``Chisel documentation,''
\url{https://www.chisel-lang.org/docs}, accessed Aug.~31, 2026.

\bibitem{circt}
LLVM Project, ``CIRCT documentation,''
\url{https://circt.llvm.org/docs/}, accessed Sep.~1, 2026.
\end{thebibliography}

\end{document}
